\section{単回帰分析}
\label{s09_Regression}

\subsection{授業内容}
\subsubsection{科目の中でのこのコマの位置づけ}
\begin{tabular}[t]{ccccccc}
    記述統計[3] & 確率[1] & モデル[3] & 推測・推定[0] & 検定[0] & 実践[2]
\end{tabular}

\subsubsection{概要}
古典的テスト理論から測定に含まれる誤差を知り，誤差が正規分布に従っていると考えることで，平均の位置を推定することができるようになった。ここでは他の変数との関係から，平均の位置が関数で表現される最も単純な場合，線形モデルについて考える。
\subsubsection{コマ主題細目(3-5項目箇条書き)}
\begin{description}
    \item[相関係数とモデル]記述統計の中で学んだ複数の変数間関係を記述する共分散，相関係数について再確認する。とくに相関係数が線形関係の指標になっていたことに留意する。しかし逆に考えれば，線形関係が見出せるような散布図であれば，我々が「見出している」関数がどのような形であるかを考えることができる。
    \begin{flushright}
        $\to$ \citet{kosugi2018}のP.51--52.
    \end{flushright}

    \item[予測と線形モデル]変数間関係を記述する最も単純な形として，一次関数を用いたものを考える。現象に対して関数で記述する数学的現象主義，変数を説明変数と被説明変数として捉えること，予測に誤差が加わって現象となることなど，これまでの議論をデータと数式に置き換えただけであることをしっかりと理解する。さらに単回帰分析の数式的記述について理解する。とくにデータ($X_i,Y_i$)，予測値($\hat{Y}$)，誤差($e_i$)，回帰係数($b_0,b_1$)として表される記号の細部にまで注意を払い，その意味を理解する。
    \begin{flushright}
        $\to$ \citet{豊田201706}のP.27--28.
    \end{flushright}

    \item[最小二乗法]最小二乗法によって推定値が算出できることを理解する。数式の細かい展開はここでは行わず，最小化する目的関数の直観的理解に止める。また算出された係数(傾きと切片)がどのような意味を持つのか改めて理解する。
    \begin{flushright}
        $\to$　\citet{豊田201706}のP.30
    \end{flushright}

    \item[残差と決定係数]回帰係数が計算されれば，実際のデータに対して予測値を計算することができ，予測がどの程度適合していたかを検証することができる。予測値と被説明変数との差分を残差residualsと呼ぶが，この両者の相関係数を重相関係数，その二乗したものを決定係数と呼んでモデルの適合度を測る指標にすることを理解する。
    \begin{flushright}
        $\to$\citet{野島201903}のP.58-60.
    \end{flushright}

\end{description}
\subsubsection{キーワード}
\begin{itemize}
    \item 一次関数
    \item 予測値，誤差，回帰係数
    \item 最小二乗法
\end{itemize}
\subsection{授業情報}
\paragraph{コマの展開方法}
講義
\paragraph{標準シラバスにおける位置づけ}
科目番号5 心理学統計法；1.心理学で用いられる統計手法；C 記述統計：回帰
\subsubsection{予習・復習課題}
\paragraph{予習}
共分散および相関係数の算出方法について，第3回の講義を参考に復習しておくこと。
\paragraph{復習}
\citet{山内201001}のP.62--81は相関係数と回帰分析を1つの章にまとめて解説しているので，予習と復習の両方に適した参考書である。この本は続く章に母集団と平均の言及がある，すなわち回帰分析の最小二乗法による解の段階では，記述統計の範囲で議論できることではあることにも気づかされる。最小二乗法による回帰係数の算出については，偏微分連立方程式を用いれば簡単に計算できるが，偏微分連立方程式とは何かに遡って解説している本当して\citet{kosugi2018}のP.104--112がある。また最小二乗法に頼らずに算出することも，手数はかかるが不可能ではない。その場合については\citet{西内201712}を参考のこと。概念的理解が困難に思える場合は，\citet{kosugi2018}のP.54--57を参考に実際にRで計算して確かめてみると良い。
